\documentclass[10pt,a4paper]{amsart}
\usepackage{amsmath,amssymb,amsthm,graphicx}
\usepackage{verbatim}
\newenvironment{reference}{\comment}{\endcomment}
%\newenvironment{reference}{}{}
\newenvironment{slogan}{\comment}{\endcomment}
\newenvironment{history}{\comment}{\endcomment}

% For commutative diagrams we use Xy-pic
\usepackage[all]{xy}

% We use 2cell for 2-commutative diagrams.
\xyoption{2cell}
% Theorem environments.
%
\theoremstyle{plain}
\newtheorem{theorem}[subsection]{Theorem}
\newtheorem{proposition}[subsection]{Proposition}
\newtheorem{lemma}[subsection]{Lemma}

\theoremstyle{definition}
\newtheorem{definition}[subsection]{Definition}
\newtheorem{example}[subsection]{Example}
\newtheorem{exercise}[subsection]{Exercise}
\newtheorem{situation}[subsection]{Situation}

\theoremstyle{remark}
\newtheorem{remark}[subsection]{Remark}
\newtheorem{remarks}[subsection]{Remarks}

\numberwithin{equation}{subsection}

% Macros
%
\def\lim{\mathop{\mathrm{lim}}\nolimits}
\def\colim{\mathop{\mathrm{colim}}\nolimits}
\def\Spec{\mathop{\mathrm{Spec}}}
\def\Hom{\mathop{\mathrm{Hom}}\nolimits}
\def\Ext{\mathop{\mathrm{Ext}}\nolimits}
\def\SheafHom{\mathop{\mathcal{H}\!\mathit{om}}\nolimits}
\def\SheafExt{\mathop{\mathcal{E}\!\mathit{xt}}\nolimits}
\def\Sch{\mathit{Sch}}
\def\Mor{\mathop{\mathrm{Mor}}\nolimits}
\def\Ob{\mathop{\mathrm{Ob}}\nolimits}
\def\Sh{\mathop{\mathit{Sh}}\nolimits}
\def\NL{\mathop{N\!L}\nolimits}
\def\CH{\mathop{\mathrm{CH}}\nolimits}
\def\proetale{{pro\text{-}\acute{e}tale}}
\def\etale{{\acute{e}tale}}
\def\QCoh{\mathit{QCoh}}
\def\Ker{\mathop{\mathrm{Ker}}}
\def\Im{\mathop{\mathrm{Im}}}
\def\Coker{\mathop{\mathrm{Coker}}}
\def\Coim{\mathop{\mathrm{Coim}}}

% Boxtimes
%
\DeclareMathSymbol{\boxtimes}{\mathbin}{AMSa}{"02}

%
% Macros for moduli stacks/spaces
%
\def\QCohstack{\mathcal{QC}\!\mathit{oh}}
\def\Cohstack{\mathcal{C}\!\mathit{oh}}
\def\Spacesstack{\mathcal{S}\!\mathit{paces}}
\def\Quotfunctor{\mathrm{Quot}}
\def\Hilbfunctor{\mathrm{Hilb}}
\def\Curvesstack{\mathcal{C}\!\mathit{urves}}
\def\Polarizedstack{\mathcal{P}\!\mathit{olarized}}
\def\Complexesstack{\mathcal{C}\!\mathit{omplexes}}
% \Pic is the operator that assigns to X its picard group, usage \Pic(X)
% \Picardstack_{X/B} denotes the Picard stack of X over B
% \Picardfunctor_{X/B} denotes the Picard functor of X over B
\def\Pic{\mathop{\mathrm{Pic}}\nolimits}
\def\Picardstack{\mathcal{P}\!\mathit{ic}}
\def\Picardfunctor{\mathrm{Pic}}
\def\Deformationcategory{\mathcal{D}\!\mathit{ef}}
\usepackage{cleveref}
\crefname{originalreference}{original reference}{original references}
\Crefname{originalreference}{Original reference}{Original references}
\makeatletter
\newlabel{situation-shriek}{{0F42}{source}{}{sourcepointer-1}{}}
\newlabel{situation-shriek@cref}{{[originalreference][0][]0F42}{[1][0][]source}{}{}{}}
\makeatother
\crefname{originaltheorem}{original theorem}{original theorems}
\Crefname{originaltheorem}{Original theorem}{Original theorems}
\makeatletter
\newlabel{flat-theorem-nagata}{{0F41}{source}{}{sourcepointer-2}{}}
\newlabel{flat-theorem-nagata@cref}{{[originaltheorem][0][]0F41}{[1][0][]source}{}{}{}}
\makeatother
\crefname{originallemma}{original lemma}{original lemmas}
\Crefname{originallemma}{Original lemma}{Original lemmas}
\makeatletter
\newlabel{flat-lemma-compactifyable}{{0A9Z}{source}{}{sourcepointer-3}{}}
\newlabel{flat-lemma-compactifyable@cref}{{[originallemma][0][]0A9Z}{[1][0][]source}{}{}{}}
\makeatother
\crefname{originallemma}{original lemma}{original lemmas}
\Crefname{originallemma}{Original lemma}{Original lemmas}
\makeatletter
\newlabel{lemma-twisted-inverse-image}{{0A9E}{source}{}{sourcepointer-4}{}}
\newlabel{lemma-twisted-inverse-image@cref}{{[originallemma][0][]0A9E}{[1][0][]source}{}{}{}}
\makeatother
\crefname{originallemma}{original lemma}{original lemmas}
\Crefname{originallemma}{Original lemma}{Original lemmas}
\makeatletter
\newlabel{cohomology-lemma-internal-hom}{{08DJ}{source}{}{sourcepointer-5}{}}
\newlabel{cohomology-lemma-internal-hom@cref}{{[originallemma][0][]08DJ}{[1][0][]source}{}{}{}}
\makeatother
\crefname{originallemma}{original lemma}{original lemmas}
\Crefname{originallemma}{Original lemma}{Original lemmas}
\makeatletter
\newlabel{perfect-lemma-cohomology-base-change}{{08EU}{source}{}{sourcepointer-6}{}}
\newlabel{perfect-lemma-cohomology-base-change@cref}{{[originallemma][0][]08EU}{[1][0][]source}{}{}{}}
\makeatother
\crefname{originallemma}{original lemma}{original lemmas}
\Crefname{originallemma}{Original lemma}{Original lemmas}
\makeatletter
\newlabel{lemma-dualizing-schemes}{{0A89}{source}{}{sourcepointer-7}{}}
\newlabel{lemma-dualizing-schemes@cref}{{[originallemma][0][]0A89}{[1][0][]source}{}{}{}}
\makeatother
\title{The upper shriek is computed by the duality involution applied to the pullback}
\author{The Stacks Project Authors}
\date{stacks-project:duality.tex:local-snapshot-20261012; GNU FDL 1.2 or later}
\begin{document}
\maketitle
\noindent\textit{Editorial note.} All delivered mathematical text is the author's own, copied from the pinned Stacks Project chapter file duality.tex; no statement or proof step was rewritten, reordered or removed.\par\medskip
\noindent\textit{Editorial note.} Preamble provenance: the pinned chapter file opens with an input of the shared preamble and carries no preamble text of its own, so the project's single Stacks-wide preamble staged at sources/stacks-duality/preamble.tex is excerpted verbatim under the recipe fixed by LANE\_BRIEF 22.1, as three spans, lines 8-13 (verbatim package and the reference/slogan/history environments), lines 14-18 (Xy-pic setup) and lines 154-217 (the native theorem declarations and mathematical macros). Nothing was added, removed or redefined; the wrapper class and the amsmath/amssymb/amsthm/graphicx base packages are supplied by the assembler.\par\medskip
\noindent\textit{Editorial note.} The counted statement and proof contain no cross reference, no citation and no included file, so the delivered body is self contained.\par\medskip
\noindent\textit{Editorial note.} Canonical location of the counted proposition: Stacks tag 0AU2 (sources/stacks-duality/pinned-original-source.tex line 4734).\par\medskip
\noindent\textit{Allowed standard prerequisites (editorial).}\begin{itemize}
\item The right adjoint a of pushforward along a compactification and the description f\textasciicircum{}! = j\textasciicircum{}* o a. Same chapter, Lemma labelled lemma-twisted-inverse-image, cited at proof lines 4750-4754. Reduces the statement to the compactified case.
\item The duality involution on the bounded coherent derived category of a Noetherian scheme with a dualizing complex. Same chapter, Lemma labelled lemma-dualizing-schemes, cited at proof lines 4785-4787. Gives D\_Y(D\_Y(M)) = M, the key step of the chain.
\item The internal Hom adjunction and its behaviour under the tensor product with a pullback. Cohomology, Lemma labelled cohomology-lemma-internal-hom, cited at proof lines 4764-4766, 4776-4778 and 4791-4793. Turns the Hom groups into tensor products and back.
\item Projection formula style base change for the cohomology of a tensor product. Derived Categories of Schemes, Lemma labelled perfect-lemma-cohomology-base-change, cited at proof lines 4770-4772. Moves the pullback out of the pushforward.
\item Nagata compactification. More on Flatness, Theorem labelled flat-theorem-nagata and Lemma labelled flat-lemma-compactifyable, cited at proof lines 4748-4750. Provides the compactification used throughout.
\end{itemize}\medskip
\ifdefined\hypertarget\hypertarget{sourcepointer-1}{}\fi
\noindent\textit{Source reference.} Original reference 0F42: sources/stacks-duality/pinned-original-source.tex, source lines 3953--3953. Same-chapter pointer into the author environment situation at source line 3953; invoked at proof line None of the counted proof. It is retained only as a navigation pointer; every load-bearing step is inside the delivered span.\par\medskip
\ifdefined\hypertarget\hypertarget{sourcepointer-2}{}\fi
\noindent\textit{Source reference.} Original theorem 0F41: sources/stacks-chow/tags.txt, source lines 18886--18886. Cross-chapter standard prerequisite located in the project tag list staged in the local snapshot; invoked at proof line 4750 of the counted proof. It is retained only as a navigation pointer; every load-bearing step is inside the delivered span.\par\medskip
\ifdefined\hypertarget\hypertarget{sourcepointer-3}{}\fi
\noindent\textit{Source reference.} Original lemma 0A9Z: sources/stacks-chow/tags.txt, source lines 12967--12967. Cross-chapter standard prerequisite located in the project tag list staged in the local snapshot; invoked at proof line 4751 of the counted proof. It is retained only as a navigation pointer; every load-bearing step is inside the delivered span.\par\medskip
\ifdefined\hypertarget\hypertarget{sourcepointer-4}{}\fi
\noindent\textit{Source reference.} Original lemma 0A9E: sources/stacks-duality/pinned-original-source.tex, source lines 653--653. Same-chapter pointer into the author environment lemma at source line 653; invoked at proof line 4752 of the counted proof. It is retained only as a navigation pointer; every load-bearing step is inside the delivered span.\par\medskip
\ifdefined\hypertarget\hypertarget{sourcepointer-5}{}\fi
\noindent\textit{Source reference.} Original lemma 08DJ: sources/stacks-chow/tags.txt, source lines 10639--10639. Cross-chapter standard prerequisite located in the project tag list staged in the local snapshot; invoked at proof line 4785 of the counted proof. It is retained only as a navigation pointer; every load-bearing step is inside the delivered span.\par\medskip
\ifdefined\hypertarget\hypertarget{sourcepointer-6}{}\fi
\noindent\textit{Source reference.} Original lemma 08EU: sources/stacks-chow/tags.txt, source lines 10685--10685. Cross-chapter standard prerequisite located in the project tag list staged in the local snapshot; invoked at proof line 4788 of the counted proof. It is retained only as a navigation pointer; every load-bearing step is inside the delivered span.\par\medskip
\ifdefined\hypertarget\hypertarget{sourcepointer-7}{}\fi
\noindent\textit{Source reference.} Original lemma 0A89: sources/stacks-duality/pinned-original-source.tex, source lines 472--472. Same-chapter pointer into the author environment lemma at source line 472; invoked at proof line 4791 of the counted proof. It is retained only as a navigation pointer; every load-bearing step is inside the delivered span.\par\medskip
\begin{lemma}
\label{lemma-shriek-via-duality}
In Situation \ref{situation-shriek} let $f : X \to Y$ be a morphism of
$\textit{FTS}_S$. Let $K$ be a dualizing complex
on $Y$. Set $D_Y(M) = R\SheafHom_{\mathcal{O}_Y}(M, K)$ for
$M \in D_{\textit{Coh}}(\mathcal{O}_Y)$ and
$D_X(E) = R\SheafHom_{\mathcal{O}_X}(E, f^!K)$ for
$E \in D_{\textit{Coh}}(\mathcal{O}_X)$. Then there is a canonical
isomorphism
$$
f^!M \longrightarrow D_X(Lf^*D_Y(M))
$$
for $M \in D_{\textit{Coh}}^+(\mathcal{O}_Y)$.
\end{lemma}

\begin{proof}
Choose compactification $j : X \subset \overline{X}$ of $X$ over $Y$
(More on Flatness, Theorem \ref{flat-theorem-nagata} and
Lemma \ref{flat-lemma-compactifyable}). Let $a$ be the
right adjoint of Lemma \ref{lemma-twisted-inverse-image} for
$\overline{X} \to Y$. Set
$D_{\overline{X}}(E) = R\SheafHom_{\mathcal{O}_{\overline{X}}}(E, a(K))$
for $E \in D_{\textit{Coh}}(\mathcal{O}_{\overline{X}})$.
Since formation of $R\SheafHom$ commutes with restriction to opens
and since $f^! = j^* \circ a$ we see that it suffices to prove that
there is a canonical isomorphism
$$
a(M) \longrightarrow D_{\overline{X}}(L\overline{f}^*D_Y(M))
$$
for $M \in D_{\textit{Coh}}(\mathcal{O}_Y)$. For
$F \in D_\QCoh(\mathcal{O}_X)$ we have
\begin{align*}
\Hom_{\overline{X}}(
F, D_{\overline{X}}(L\overline{f}^*D_Y(M)))
& =
\Hom_{\overline{X}}(
F \otimes_{\mathcal{O}_X}^\mathbf{L} L\overline{f}^*D_Y(M), a(K)) \\
& =
\Hom_Y(
R\overline{f}_*(F \otimes_{\mathcal{O}_X}^\mathbf{L} L\overline{f}^*D_Y(M)),
K) \\
& =
\Hom_Y(
R\overline{f}_*(F) \otimes_{\mathcal{O}_Y}^\mathbf{L} D_Y(M),
K) \\
& =
\Hom_Y(
R\overline{f}_*(F), D_Y(D_Y(M))) \\
& =
\Hom_Y(R\overline{f}_*(F), M) \\
& = \Hom_{\overline{X}}(F, a(M))
\end{align*}
The first equality by Cohomology, Lemma \ref{cohomology-lemma-internal-hom}.
The second by definition of $a$.
The third by Derived Categories of Schemes, Lemma
\ref{perfect-lemma-cohomology-base-change}.
The fourth equality by Cohomology, Lemma \ref{cohomology-lemma-internal-hom}
and the definition of $D_Y$.
The fifth equality by Lemma \ref{lemma-dualizing-schemes}.
The final equality by definition of $a$.
Hence we see that $a(M) = D_{\overline{X}}(L\overline{f}^*D_Y(M))$
by Yoneda's lemma.
\end{proof}

\end{document}
